\chapter{逻辑回归与概率分类}
\chaptergoal{能够解释分类预测与连续目标预测的区别；理解线性得分、Sigmoid、Softmax和交叉熵的作用；使用 sklearn 完成二分类和多分类；通过阈值变化分析精确率与召回率；把逻辑回归作为后续经典模型和 MLP 的基线。}

逻辑回归（logistic regression）虽然含有“回归”二字，却主要用于分类。它给输入样本计算一个得分，再把得分转换为类别概率。其意义不仅在于实现简单，更在于把线性模型、概率预测、损失函数和优化过程联系起来。第12章中我们将看到：在这种结构上加入隐藏层，就可以进入前馈神经网络。
\section{从连续目标到类别目标}
前一章回归的目标是连续数值，比如声压级。分类则要求模型回答一条样本属于哪一个类别。对于 Iris，目标可取三种花的类别；对天体目录，目标可为恒星、星系或类星体。用普通线性回归拟合 $0,1,2$ 等类别编码会人为引入类别之间的数值距离。分类器应根据类别性质设计预测输出和训练目标。

先考虑二分类。用 $y\in\{0,1\}$ 表示类别，模型输出属于类别1的概率估计 $p(\mathbf x)$。只有在阈值确定之后，概率才被转换成最终类别。训练过程中使用连续的概率输出，使算法有机会区分“接近边界”和“非常有把握”的预测；但经过有限数据训练得到的概率并非自动经过良好校准，实际应用仍需检验。

\subsection{为什么不能直接把类别代码当普通数值回归}
对二分类目标$y\in\{0,1\}$，可以把线性回归输出当作一个得分，但直线本身没有把输出限制在0与1之间。远离训练区域时，它可能给出$-0.4$或1.3，难以直接解释为概率。更重要的是，分类任务关注类别及相应决策代价，不能只因为代码存为数字就忽略目标含义。

逻辑回归选择先计算任意实数得分，再把得分映射为概率。它仍然可以借助第3章的损失与优化思路训练，但模型输出的解释与线性回归不同。名称中的“回归”源于对概率的某种变换建立线性关系，不意味着本章在预测连续物理量。

\subsection{概率估计与最终行动分开}
假设一个天体候选被模型赋予0.7的目标类别概率。对“优先安排后续观测”与“写入高纯度目录”，可能采用不同的阈值，因此同一个概率可以导致不同动作。这不是模型自相矛盾，而是应用代价不同。

在学习阶段先把两层分开：概率模型提供对类别的估计，决策规则把估计转成行动。验证时也分别检查概率质量和阈值下的错误。后面会看到，准确率相近的两个模型，交叉熵和概率校准可能明显不同。

\section{线性得分与 Sigmoid 概率}

\textbf{得分}是模型计算出的中间数值，可以为负、为零或为正，它本身还不是概率。\textbf{概率估计}则是模型对某一类别支持程度的数值表达，二分类时用0到1之间的数表示。逻辑回归把两步分开：先求线性得分，再把得分变成概率。

考虑人工给定的输入 $\mathbf x=(2,0.5)$、权重 $\mathbf w=(1,-1)$ 和截距 $b=-0.5$。得分为 $z=1\times2-1\times0.5-0.5=1$。这只是检查预测计算的例子，权重尚不是从某个真实数据集学出的结果。

首先计算线性得分 $z=\mathbf w^{\mathsf T}\mathbf x+b$，再使用
\begin{equation}
\sigma(z)=\frac{1}{1+\exp(-z)},\quad 0<\sigma(z)<1,
\end{equation}
把实数变换到 $(0,1)$。当 $z=0$ 时，输出为 $0.5$；$z$ 越大，输出越接近1。这并不是对类别编号做了线性拟合，而是用线性得分构造类别概率模型。

在刚才的例子中，$p=1/(1+e^{-1})\approx0.7311$。若规定 $p\geq0.5$ 判为类别1，就预测为1；若应用要求 $p\geq0.8$ 才判为1，同一个概率将得到类别0。前一种变化是改变判断规则，不是重新计算模型参数。

\textbf{决策边界}指模型从预测一个类别切换为另一个类别的位置集合。对于0.5阈值，分界处满足 $z=0$。例如二维得分 $x_1-x_2-0.5=0$ 对应一条直线，它把平面分成两个预测区域。“超平面”只是直线和平面在更高维坐标中的推广名称。


\fig{ch06_sigmoid.pdf}{Sigmoid 将线性得分映射到0与1之间。二分类中，阈值0.5对应得分0；应用阈值可依任务改变。}

在二维输入时，$z=0$ 是一条直线；在更多维时是一张超平面。即使输出经过非线性 Sigmoid 映射，逻辑回归在原始输入空间里的二分类边界仍是线性的。它不能凭增加迭代次数自动学习任意弯曲的边界；要表达更复杂的关系，可以构造非线性特征或更换模型。

\subsection{从概率的对数几率得到Sigmoid}
对正类概率$p\in(0,1)$，\textbf{几率}定义为$p/(1-p)$，表示正类与负类概率的比。$p=0.8$时几率为4，即正类概率是负类的四倍；它不是概率4。几率始终为正，而它的对数可以取任意实数。

逻辑回归假定对数几率由输入的线性组合描述：
\[
\log\frac{p}{1-p}=z=b+\mathbf x^{\mathsf T}\mathbf w.
\]
两边取指数并解出$p$，有$p=e^z(1-p)$，因此$p=e^z/(1+e^z)=1/(1+e^{-z})$。这就得到Sigmoid映射，而不是在直线后面任意接一条弯曲函数。这个假设并不要求输入变量服从正态分布。

在其他输入保持不变的数学条件下，$x_j$增加一个单位会使几率乘以$e^{w_j}$。这是模型内部的条件关联解释，不自动成为因果效应。若先标准化输入，一个单位指一个训练标准差，解释时仍需换回原尺度或明确说明。

\worked{阈值改变怎样移动边界}
一维得分为$z=2x-1$。当阈值$p\ge0.5$时，等价于$z\ge0$，即$x\ge0.5$。若阈值提高到0.8，需$z\ge\log(0.8/0.2)=\log4\approx1.3863$，边界变成$x\ge1.1931$。系数没有重新训练，变化的是判断规则。

在多维二分类中，固定阈值对应$\mathbf x^{\mathsf T}\mathbf w+b=\log(t/(1-t))$，仍然是一张超平面；提高阈值主要移动这张平面，而不会使边界自动弯曲。若问题确实需要曲线边界，可以加入非线性特征或选择后续的核方法、树和神经网络。

\section{交叉熵损失与正则化}

模型给出了类别概率以后，还需要衡量它是否支持了已知目标。\textbf{交叉熵损失}在本章指下面这个比较已知类别与预测概率的函数。先不要求记住其名称，而是观察同一个真实类别下，模型给出高概率和低概率分别付出多大代价。公式中的 $\log$ 指自然对数。

对二分类标签 $y_i\in\{0,1\}$ 和预测概率 $p_i$，常见单样本交叉熵损失为
\begin{equation}
\ell_i=-y_i\log p_i-(1-y_i)\log(1-p_i).
\end{equation}
若真实类别为1，模型给它的概率越小，惩罚就越大；真实类别为0时情况相反。

例如真实类别 $y=1$，预测概率 $p=0.8$ 时，损失为 $-\log0.8\approx0.2231$；若 $p=0.2$，损失变成 $-\log0.2\approx1.6094$。二者都可能在不同阈值下形成类别输出，但损失明确区分它们对真实类别的支持程度。若真实类别为0，对应的是 $-\log(1-p)$，上述两个概率的损失大小交换。
这样，训练不仅要求类别最终正确，还鼓励概率对已知类别的支持程度合乎目标函数。实践中概率接近0或1的计算需要数值稳定处理，成熟库内部会采用相应的实现，不建议在教学练习中直接逐样本取对数而不防溢出。

正则化能抑制过大的权重并影响边界。sklearn \texttt{LogisticRegression} 常使用参数 \texttt{C} 控制正则化强度，其直觉是：\texttt{C} 较小通常意味着更强的正则化。不同求解器对惩罚类型、类别问题和收敛特性的支持可能不同；使用接口时应查阅对应软件版本文档。改变 \texttt{C} 后应依据验证结果比较，而不是单看训练准确率。

\subsection{从已知类别的概率到训练目标}
对于一条标签$y\in\{0,1\}$的记录，模型赋予已知类别的概率可以统一写成$p^y(1-p)^{1-y}$。当$y=1$时，它就是$p$；当$y=0$时，它就是$1-p$。若暂按给定输入后记录独立处理，整批标签的概率由这些项相乘。

希望已知标签得到较大概率，等价于增大这些项的乘积。对乘积取自然对数，把乘法变成求和；再取负号，得到可以最小化的负对数似然：
\[
J(\mathbf w,b)=-\frac1n\sum_i\{y_i\log p_i+(1-y_i)\log(1-p_i)\}.
\]
这就是本章使用的平均二元交叉熵。只需这几步就能解释公式从哪里来，无需在入门时展开一般极大似然理论。概率模型和记录关系的假设仍应根据应用检查，尤其是重复对象不能凭这条公式变成独立样本。

\subsection{概率严重错误为什么受到更大惩罚}
真实标签为1时，给$p=0.9$的损失约0.105，给$p=0.1$约2.303，给$p=0.001$约6.908。三个结果不只是“对或错”的差别，还表示模型对错误结论有多强的支持。相同准确率下，一个把错误样本赋予极端概率的模型，可能具有更差的交叉熵。

Sigmoid的输出在数学上位于开区间$(0,1)$，计算机中却可能因有限精度变得极接近边界。因此成熟实现通常直接使用稳定的得分形式计算损失，而不是先得到被舍入为0的概率再取对数。教学手算可以使用普通公式，实际程序应采用经过测试的损失函数或稳定的\texttt{logaddexp}运算。

\subsection{交叉熵的梯度怎样连接第3章}
Sigmoid导数为$dp/dz=p(1-p)$。对单样本交叉熵求导，链式法则中的分母与该导数相消，得到$\partial\ell/\partial z=p-y$。于是
\[
\nabla_{\mathbf w}J=\frac1nX^{\mathsf T}(\mathbf p-\mathbf y),\qquad
\frac{\partial J}{\partial b}=\frac1n\sum_i(p_i-y_i).
\]
这条式子与线性回归梯度的结构很相似，但残差是概率与0/1目标之差。梯度下降仍是“计算方向，再按学习率更新”，只不过模型预测与目标函数已经改变。

\worked{手算一次概率模型训练}
取人工记录$x=(-1,0,1)$、$y=(0,0,1)$，初始$w=b=0$，三条概率都是0.5。梯度为$\partial J/\partial w=-1/3$、$\partial J/\partial b=1/6$。取学习率0.3，同时更新得到$w=0.1,b=-0.05$。得分变为$(-0.15,-0.05,0.05)$，平均交叉熵从$\log2\approx0.693147$降至约0.652625。随书脚本按这些数字直接复算。

这一步并不保证任意学习率都能改善目标。也不能因为三条训练记录上损失降低，就直接宣称在新样本上更可靠。训练机制在这里与第3章一致，评价职责仍属于第4章。

\subsection{为什么要限制系数过大}
若训练类别可以由一条直线完全分开，无惩罚逻辑回归可能通过不断增大系数，把训练中每个正确类别的概率推得更接近1。损失继续改善，却未必得到有限、稳定的系数。正则化可以约束这一过程，并减少有限样本中极端概率的风险。

本书使用的\texttt{LogisticRegression}以\texttt{C}表达正则化强度的倒数含义，较小的\texttt{C}通常对应较强惩罚。准确数值还涉及实现中的损失归一化、样本权重和求解器，不要把它直接等同于Ridge的\texttt{alpha}。模型比较时固定特征和验证方案，再改变少量有解释意义的配置。

\section{使用 sklearn 完成二分类}
从 Iris 中选择两种花，使用其中两个测量特征展示二维决策边界，再用全部四个特征完成较正式的评价。二维图便于理解，但它与全特征模型的结果不能混为一谈。

\begin{lstlisting}[caption={鸢尾花两类问题的可复现逻辑回归示例}]
from sklearn.datasets import load_iris
from sklearn.model_selection import train_test_split
from sklearn.preprocessing import StandardScaler
from sklearn.linear_model import LogisticRegression
from sklearn.pipeline import make_pipeline
from sklearn.metrics import classification_report

X, y = load_iris(return_X_y=True)
mask = y < 2
X, y = X[mask], y[mask]
X_train, X_val, y_train, y_val = train_test_split(
    X, y, stratify=y, test_size=0.25, random_state=42
)
model = make_pipeline(StandardScaler(), LogisticRegression(max_iter=2000))
model.fit(X_train, y_train)
prob = model.predict_proba(X_val)[:, 1]
pred = model.predict(X_val)
print(classification_report(y_val, pred))
\end{lstlisting}

\texttt{predict()} 返回类别，\texttt{predict\_proba()} 返回各类别的概率估计。不能把 \texttt{predict\_proba()} 的第二列不加检查地解释成任意指定科学类别：要通过模型的 \texttt{classes\_} 属性核对类别顺序。

\fig{ch06_logistic_boundary.pdf}{Iris 两类、两个特征上的逻辑回归边界示意。正式全特征评价不能用二维演示图代替。}

\subsection{逐个核对接口返回的对象}
\texttt{fit(X,y)}估计系数；\texttt{decision\_function(X)}返回类别得分；\texttt{predict\_proba(X)}返回类别概率；\texttt{predict(X)}给出最终类别。对于二分类，手工阈值处理必须先核对\texttt{classes\_}，保证取出的那一列确实对应所关心的正类。

如果模型放在流水线中，系数属于已经变换后的特征空间。最后一个模型的\texttt{coef\_}和\texttt{intercept\_}不能直接与原始厘米数代入，除非同时应用前面的标准化。可以用几条输入比较流水线预测与手工“缩放—线性得分—Sigmoid”计算，误差应在数值精度范围内。这个检查比只看最终准确率更能发现列序和变换遗漏。

\subsection{简单二分类与困难二分类分别承担什么用途}
鸢尾花中选择一个较易区分的物种，与其他物种比较，适合第一次认识概率与边界；选择测量分布更重叠的类别，更适合观察阈值取舍。不能要求每个数据划分都必然出现某种教科书曲线。若三个阈值给出完全相同预测，应检查概率是否都远离这些阈值，而不是伪造分数变化。

本章采用两种证据：固定的人工概率列表保证可以逐条看见阈值怎样改变结果；真实标准数据实验检查模型是否确实给出相应行为。前者解释计算，后者展示在具体数据上的表现。图表明确标注二者来源。

\section{多分类与 Softmax 的直观认识}

\textbf{多分类}指每个样本需要从三个或更多互斥类别中选一个；不是同时给样本贴上多个可共存的标签。葡萄酒化学分析案例（第\ref{sec:wine-intro}节）有三个类别，适合说明这种任务。\textbf{Softmax}是把一组实数得分转成总和为1的一组正数的运算，下面用这组数表示各类别概率。

对于三类或更多类别，可为每类计算一个得分 $z_k$，再利用 Softmax 使输出非负且和为1：
\begin{equation}
\hat p_k=\frac{\exp(z_k)}{\sum_j\exp(z_j)}.
\end{equation}
在数值实现中，可先减去所有得分的最大值，以避免计算指数时出现溢出。

例如三个类别的得分分别为 $\log2,\log3,0$，指数分别为2、3、1，总和为6，三个概率就是 $1/3,1/2,1/6$。按最大概率预测时选择第二类。概率之和为1，但最大值只有0.5；多分类中并不要求最大概率大于0.5才能产生类别预测。
多分类逻辑回归仍然主要通过线性得分形成分类边界。Wine 有13个输入特征与3种类别，适合在更多特征下研究标准化、交叉验证与系数解释。

\begin{lstlisting}[caption={葡萄酒类别的逻辑回归与交叉验证}]
from sklearn.datasets import load_wine
from sklearn.model_selection import StratifiedKFold, cross_val_score
from sklearn.pipeline import make_pipeline
from sklearn.preprocessing import StandardScaler
from sklearn.linear_model import LogisticRegression

X, y = load_wine(return_X_y=True)
from sklearn.model_selection import train_test_split
X, X_hold, y, y_hold = train_test_split(
    X, y, test_size=0.2, stratify=y, random_state=42
)
# All comparisons below use development X, y only.
clf = make_pipeline(StandardScaler(), LogisticRegression(max_iter=3000))
cv = StratifiedKFold(n_splits=5, shuffle=True, random_state=7)
score = cross_val_score(clf, X, y, cv=cv, scoring="accuracy")
print(score.mean(), score.std())
\end{lstlisting}

这里先保留20\%的数据，再只在其余开发部分进行交叉验证。代码中的\texttt{X, y}此时指开发数据；\texttt{X\_hold, y\_hold}不参与本章的模型选择。

\subsection{多分类交叉熵怎样选择概率向量}
对$K$个互斥类别，模型输出$K$个概率，和为1。若第$i$条样本参考类别为$c_i$，单样本损失为$-\log p_{i,c_i}$。也可以把目标写成独热向量$y_{ik}$，得到$-\sum_k y_{ik}\log p_{ik}$。这里只取已知类别对应的对数，但由于Softmax把所有类别耦合，调整一个得分会改变整个概率向量。

例如真实类别为第2类，两个模型分别给$(0.4,0.5,0.1)$和$(0.05,0.9,0.05)$，都预测第2类，准确率上没有区别，损失却分别为$-\log0.5$和$-\log0.9$。若后续任务依赖候选排序或概率决策，就需要了解这种差别。反过来，概率很尖锐也不自动代表校准良好，错误地自信会受到更大损失。

\subsection{一对其余分类与联合多分类的区别}
一种思路是为每个类别单独训练“该类对其余类”的二分类器；另一种思路是用Softmax联合训练全部类别。前者的各个二分类输出不天然和为1，后者在概率模型中直接施加总和约束。两者都能用于多分类，但训练目标和概率解释不完全相同。

本书主线使用软件明确支持的联合逻辑回归配置，不要求学生遍历所有实现。阅读他人程序时，仍应查看模型究竟采用哪一种策略。仅看到三个输出列，不能推断其已经满足同一个概率模型。

\section{概率、阈值与分类错误}
二分类中，如果默认以0.5为阈值，把 $p\geq0.5$ 判为正类；但科学任务常需要不同权衡。筛选稀有天体时，较低阈值可能增加候选的覆盖率，也会引入更多误报；后续人工检查能力有限时，阈值选择还受资源约束。正确的教学问题不是“哪个阈值永远最好”，而是“在何种代价和资源条件下，选择什么阈值”。

\begin{lstlisting}[caption={通过开发集概率检查阈值的取舍}]
from sklearn.metrics import precision_score, recall_score
for threshold in [0.2, 0.5, 0.8]:
    labels = (prob >= threshold).astype(int)
    print(threshold,
          precision_score(y_val, labels, zero_division=0),
          recall_score(y_val, labels, zero_division=0))
\end{lstlisting}

这段代码承接前面的二分类演示，使用同一验证部分\texttt{X\_val/y\_val}。它用于选择阈值，因此不再把该部分称为独立测试集。正式项目还需另行预留最终测试数据。阈值不会改变模型计算的概率，却会改变输出类别和相应的混淆矩阵。

\worked{八条人工概率记录的阈值实验}
取概率$(0.95,0.80,0.60,0.55,0.40,0.30,0.10,0.05)$，对应参考类别$(1,1,0,1,0,1,0,0)$。这些数由教学构造，不是某次模型运行伪装成的输出。规定$p\ge t$报正类。

当$t=0.2$时，前6条入选，TP=4、FP=2、FN=0、TN=2，精确率$4/6$、召回率1。当$t=0.5$时，前4条入选，TP=3、FP=1、FN=1、TN=3，精确率和召回率均为0.75。当$t=0.8$时，仅前2条入选，TP=2、FP=0、FN=2、TN=4，精确率1、召回率0.5。每一项都能从列表逐条数出。

\fig{v3_ch06_thresholds.pdf}{人工概率列表与三个阈值。阈值改变选择哪些记录，不重新训练概率模型；参考标签用于统计每个选择的后果。}

这个例子展示了误报与漏检的取舍，但精确率不在所有数据上随阈值单调上升。某个高概率错误样本可能在极高阈值下单独留下，使精确率反而下降。应依据实际记录计算，而不是把常见趋势写成无条件定理。

\subsection{从错误代价推一个阈值}
在概率能够用于当前总体、且正确判断代价取零的简化条件下，报正类的期望代价为$C_{\rm FP}(1-p)$，报负类为$C_{\rm FN}p$。选择报正类需满足前者不大于后者，得到
\[
p\ge\frac{C_{\rm FP}}{C_{\rm FP}+C_{\rm FN}}.
\]
漏检代价较大时，阈值降低；误报代价较大时，阈值提高。这个推导帮助理解0.5不是所有应用的最佳选择。实际候选筛选还涉及资源预算、概率校准、类别比例变化和测量质量，不能直接把主观代价代入后就停止验证。

\subsection{实际结果：准确率与概率质量可以给出不同排序}
葡萄酒开发数据的统一五折中，标准化逻辑回归$C=0.1$的平均准确率约0.9860、平均交叉熵约0.1627；$C=1$的准确率约0.9788、交叉熵约0.0732；GaussianNB分别约0.9650和0.0903。这里较高准确率的设置并没有较低交叉熵。

这不是指标互相矛盾，而是它们评价不同方面。若目标是固定阈值下少错几个类别，可以按相应指标选择；若需要可靠概率排序，就必须另外检查概率质量与校准。所有这些数字来自开发验证，尚未用最终保留数据确认，不应写成模型的普遍排名。

\section{朴素贝叶斯比较与专业迁移}

\textbf{先验概率}是在查看当前特征之前对类别比例的描述；\textbf{似然}描述假定类别已知时，出现当前特征的相对可能性；\textbf{后验概率}是在观察特征后对类别概率的更新。\textbf{贝叶斯规则}把它们联系起来：后验与“先验乘以似然”成比例，再归一化使各类概率之和为1。

用一个明确的人工计数模型说明：类别A与B的先验为0.6、0.4。某记录的两个特征都记为“高”；假设在A类中各自为高的概率是0.8、0.5，在B类中是0.2、0.25。\textbf{条件独立近似}是指在类别已知的条件下，用两个特征概率的乘积近似其联合概率。于是A、B的未归一化权重分别为 $0.6\times0.8\times0.5=0.24$ 与 $0.4\times0.2\times0.25=0.02$，归一化后A类概率为 $0.24/0.26=12/13$。这不是说两项测量在所有样本中必定独立，而是说明模型作了什么近似。

实际葡萄酒数据的输入是连续数值。\texttt{GaussianNB}是采用高斯分布描述各类别下各连续特征的一种朴素贝叶斯实现，其中使用概率密度而不是“恰好等于某个实数”的概率。训练时根据各类记录估计均值和方差。它与上面的离散例子共享贝叶斯组合思想，但特征的概率模型不同。\textsuperscript{[1]}

朴素贝叶斯提供另一个概率分类基线。它利用贝叶斯规则，并在给定类别时作特征条件独立的近似。对于相关性强的物理测量，这个假设可能明显不成立，但模型仍可能在某些数据上提供有价值的预测。教学上以 \texttt{GaussianNB} 做一个对照，不将其与逻辑回归的概率输出及系数含义混淆。

\subsection*{天文方向：根据测光判断天体类别}
\label{sec:sdss-intro}
\textbf{斯隆数字巡天（Sloan Digital Sky Survey，SDSS）}提供大范围天空的观测目录。\textbf{测光}是在指定波段测量光的强弱，\textbf{光谱}则记录光随波长的分布。二者给出的信息细节不同。本案例设想在只有测光输入时预测天体类别，因此不能把只有获得光谱后才知道的结果也作为输入。\textsuperscript{[14,19]}

输入可选 $u,g,r,i,z$ 五个观测波段的星等，以及两个波段星等之差（称为颜色）。目标使用所选数据版本光谱产品中的恒星、星系、类星体类别，常见代码分别为STAR、GALAXY、QSO。恒星是像太阳那样的单个恒星天体，星系是由许多恒星等组成的系统，类星体是有强烈活动核心的遥远天体类型；本练习学习目录标签的预测，而不是重新定义这些物理类别。\textsuperscript{[14]}

一行应明确对应一个天体还是一次测量。需要合并重复观测时保留对象编号；划分数据时防止同一天体同时进入训练与评价。字段字典还必须记录测光版本、质量条件与标签筛选。这里给出任务设计，不预先宣称某个模型已经在真实目录上达到某一准确率。

\subsection*{物理方向：根据声呐回波区分目标}
\label{sec:sonar-intro}
\textbf{声呐}利用发射与接收声波研究目标。本例使用\textbf{声呐回波数据集（Sonar）}，包含208条记录，其中111条对应实验金属圆柱目标、97条对应岩石。每条记录的60个输入表示在不同频带、规定时间窗口内整理并归一化的回波能量，数值位于0到1之间。目标是两种实验对象的类别。\textsuperscript{[17]}

原始文件常用M与R表示两类；随书转换脚本的映射写在字段字典中。文件名有时带有“mines”，但此处不能把这些实验记录直接称为已经验证的实际水雷识别系统。数据没有给出每条记录的全部设备和环境信息，也不能随意给60列附加具体中心频率或物理单位。先比较逻辑回归与朴素贝叶斯，再在后续章节保持相同数据划分比较其他模型。

\keypoint{逻辑回归从 $\mathbf w^\mathsf T\mathbf x+b$ 得到概率，是从线性模型过渡到 MLP 的关键。后续加入隐藏层后，网络可以学习更丰富的输入表示，但训练、验证和科学解释的问题不会因此消失。}

\subsection{贝叶斯公式中的三个对象}
对类别$c$和特征$x$，贝叶斯公式把后验概率写成$P(c\mid x)\propto P(x\mid c)P(c)$。$P(c)$称为先验类别概率，表示看到该样本特征之前类别的比例信息；$P(x\mid c)$描述该类内特征的分布；$P(c\mid x)$是结合特征之后的类别概率。分母用于让各类概率总和为1。

朴素贝叶斯进一步作条件独立近似：给定类别以后，各特征的联合分布近似分解为$\prod_jP(x_j\mid c)$。这不是说现实中特征彼此绝对独立，而是所用概率模型的简化。\texttt{GaussianNB}对每个类别、每个连续特征拟合一个高斯分布，再按该近似组合。

\worked{先验怎样改变同一证据的判断}
某个特征在A类出现的概率为0.8，在B类为0.2，先验A、B各半时，该特征出现后A的概率为$0.8/(0.8+0.2)=0.8$。若A先验只有0.1、B为0.9，则后验A为$0.8\times0.1/(0.8\times0.1+0.2\times0.9)=0.3077$。同一特征证据得到不同后验，是因为原始总体比例不同。

这个人工算例与第4章稀有对象的基数效应一致。科学样本经过人为平衡后，直接输出的概率可能不再对应原始巡天总体，需要检查抽样和校准。条件独立近似也可能让相关特征被重复计入，造成过度自信；分类准确率尚可不表示概率已经可靠。

\subsection{由测光或回波到专业分类}
将已知的测光特征预测光谱参考类别时，必须区分提供标签的更昂贵信息和真正部署可用的输入；将声呐频带特征用于物理分类时，必须关注小样本波动及设备条件。专业案例首次介绍集中在第17章和数据字典，此处先复用统一的分类流程，不用模型名称代替任务说明。

逻辑回归的系数和概率为后面的神经网络提供参照。一个没有隐藏层、线性得分加Sigmoid输出的二分类结构，正是本章模型。加入隐藏层以后，变化的是输入表示的学习能力，交叉熵与梯度更新的基本联系仍然存在。

\chapterpractice{先完成 Iris 的二维概率边界演示，再使用 Wine 的全部特征比较逻辑回归与 GaussianNB，并研究正则化。专业迁移从 SDSS 与 Sonar 中选择一个任务，提交数据角色表、基线结果及阈值或错误类别分析。}
\chapterexercise
\begin{enumerate}
\item 为什么“逻辑回归”在本书中是分类方法而不是连续目标回归方法？
\item 当线性得分等于0时，Sigmoid 输出多少？改变类别阈值是否等于重新训练模型？
\item 若一组任务更关心漏检而不是误报，应该重点观察哪些指标？
\item 解释条件独立近似在物理测量中为什么可能不成立，却不必然使朴素贝叶斯完全不可用。
\item 试画出没有隐藏层的二分类逻辑回归结构图，为第12章做准备。
\end{enumerate}

\section*{深化练习与方法检验}
本组题强调解释、计算与验证；D编号在本章内独立计数。计算题应写出所用定义和中间步骤，实验题保留实际输出，包括不支持预期的结果。
\begin{enumerate}[label=D\arabic*.]
\item 从 $p=\sigma(z)$ 推出二分类交叉熵对得分的导数为 $p-y$，再说明为什么这让概率分类与第3章梯度下降连接起来。
\item 使用正文八条固定概率记录，分别在阈值0.2、0.5、0.8统计混淆矩阵。提高阈值后“准确率必须增加”是否正确？
\item 误报代价为1、漏检代价为9，正确决策代价为0。对一个经过校准且适用于当前总体的正类概率 $p$，推导选择正类的期望代价阈值。为什么仍需检验概率的适用性？
\item 给三类得分同时加上1000，Softmax概率理论上是否改变？写出稳定计算方式，并解释为什么多分类最高概率不必大于0.5。
\end{enumerate}
